\documentclass{situationist-preprint}

\title{The Experience Machine and Status Quo Bias}
\msid{ETH-05-b}
\msfamily{ETH-05}
\msdiscipline{Normative and Applied Ethics}
\mstype{Consequence}
\mskeywords{experience machine; status quo bias; hedonism; reversal test; experimental philosophy; well-being}

\begin{document}
\maketitle

\begin{abstract}
De Brigard's studies of the reversed experience machine, in which subjects are told that they are already connected, are the main empirical basis for the claim that the experience machine intuition reflects status quo bias. We examine the studies in detail and draw a consequence of ETH-05: if well-being includes contact with reality but does not consist in it, the pattern of responses across his conditions is largely what one should expect from unbiased judges. Only the neutral condition shows a residue that bias explains, and that residue concerns choices, not the comparison of lives on which ETH-05-a relies.
\end{abstract}

\section{The studies}

\citet{deBrigard2010} presented subjects with a reversal of Nozick's case \citep{Nozick1974}. Subjects were asked to imagine that they had been connected to an experience machine all along and were now offered the chance to disconnect and return to their real life. The vignette came in three versions, differing in what the subject was told about that real life: nothing (neutral), that it was the life of a prisoner in a maximum-security prison (negative), or that it was the life of a wealthy artist (positive). In none of the versions did a clear majority choose to disconnect; in the negative version a large majority chose to remain connected; and in the neutral and positive versions the choices were roughly divided. De Brigard concluded that the original reluctance to plug in reflects a preference for one's current situation, whatever it is, and not a judgment that contact with reality matters for well-being.

The inference relies on the reversal test of \citet{BostromOrd2006}: if a judgment that a change from $x$ to $y$ is bad is accompanied by a judgment that the reverse change from $y$ to $x$ is also bad, and no reason is given for thinking the status quo optimal, the pair of judgments is evidence of status quo bias. Status quo bias itself is well established in the study of decisions \citep{SamuelsonZeckhauser1988}, and it is commonly explained by loss aversion: losses loom larger than equivalent gains, and departures from the status quo are coded as losses \citep{Kahneman1991}.

\section{What the conclusion of ETH-05 predicts}

ETH-05 holds that well-being includes contact with reality (Corollary~3.3 of ETH-05-a). It does not hold that contact with reality is all of well-being, or that it outweighs every other good: the comparative verdict of Assumption~3.1 of ETH-05-a concerns experiential twins, lives alike in every experience. The vignettes do not compare experiential twins. Disconnecting means exchanging the machine's experiences for those of the real life, and the real lives described are, in two of three conditions, much worse or merely different in experience.

\begin{proposition}\label{p:negative}
Given the conclusion of ETH-05, a judge free of bias may prefer to remain connected in the negative condition.
\end{proposition}

\begin{proof}
The connected life has contact with reality at a low level and good experiences; the life in prison has contact with reality and very bad experiences. Since the subject's experiences are part of her well-being on every view, and ETH-05 adds contact with reality without making it lexically prior, the comparison can favor the connected life. Preferring it is then a correct judgment of well-being, not a bias.
\end{proof}

By the same reasoning, the positive condition is a comparison between a life with contact and the good experiences of a successful artist and a life without contact and the good experiences of the machine. ETH-05 predicts a preference for disconnecting, other things equal, and the roughly divided responses show this preference weaker than it should be. The neutral condition is where a residue of bias can be found: subjects are given no information about the real life, and, uncertain whether it is better or worse in experience, many choose to keep what they have. That is loss aversion under uncertainty, and the reversal test correctly diagnoses it.

\begin{corollary}\label{c:residue}
The evidence of status quo bias in the reversed experience machine is confined to the conditions in which the experiential quality of the real life is unknown or comparable, and there it concerns a choice to change one's situation.
\end{corollary}

The comparative verdict on which ETH-05-a relies involves no change of situation. Two lives are compared as completed wholes, neither of which is the judge's own, and there is no status quo to be biased toward. The reversal test has nothing to apply to, since there is no change whose reverse could be assessed. Corollary~\ref{c:residue} therefore leaves Assumption~3.1 of ETH-05-a untouched.

\begin{objection}
\citet{Weijers2014} constructed vignettes designed to remove status quo bias and other confounds, and found much weaker support for the anti-hedonist verdict. The comparative judgment may be equally fragile.
\end{objection}

\begin{reply}
Weijers's vignettes ask about choices, whether one's own or a stranger's, and they reduce the role of the status quo by making neither option the current one. That is a step toward the comparative form, and to the extent that his results weaken the verdict, they show that the choice-based versions are confounded in further ways, as ETH-05-a argues. They do not test the comparison of experiential twins, in which the lives are stipulated to be alike in experience, which removes the trade-off between experience and contact that drives the choice cases. An empirical study of that comparison is what would bear on ETH-05, and the predictions of Proposition~\ref{p:negative} show what such a study would have to control for.
\end{reply}

\section{Conclusion}

The reversed experience machine shows status quo bias where the real life's quality is unknown, and shows correct judgment where the real life is described as much worse. Neither result bears on the comparison of experiential twins, and the second is predicted by the view that well-being includes, but is not exhausted by, contact with reality.

\begin{thebibliography}{9}\small
\bibitem[Bostrom and Ord(2006)]{BostromOrd2006} Bostrom, N. and Ord, T. (2006). The reversal test: Eliminating status quo bias in applied ethics. \emph{Ethics}, 116(4), 656--679.
\bibitem[De Brigard(2010)]{deBrigard2010} De Brigard, F. (2010). If you like it, does it matter if it's real? \emph{Philosophical Psychology}, 23(1), 43--57.
\bibitem[Kahneman et~al.(1991)]{Kahneman1991} Kahneman, D., Knetsch, J.~L. and Thaler, R.~H. (1991). Anomalies: The endowment effect, loss aversion, and status quo bias. \emph{Journal of Economic Perspectives}, 5(1), 193--206.
\bibitem[Nozick(1974)]{Nozick1974} Nozick, R. (1974). \emph{Anarchy, State, and Utopia}. Basic Books.
\bibitem[Samuelson and Zeckhauser(1988)]{SamuelsonZeckhauser1988} Samuelson, W. and Zeckhauser, R. (1988). Status quo bias in decision making. \emph{Journal of Risk and Uncertainty}, 1(1), 7--59.
\bibitem[Weijers(2014)]{Weijers2014} Weijers, D. (2014). Nozick's experience machine is dead, long live the experience machine! \emph{Philosophical Psychology}, 27(4), 513--535.
\end{thebibliography}

\end{document}
